\textbf{Phần 1. Trắc nghiệm lựa chọn}
\setcounterpageref{ex}{0}

\begin{ex} %[4D4V18]
$1\text{ m}^2$ bằng bao nhiêu đề-xi-mét vuông?
\choice
{\True $100\text{ dm}^2$}
{$10\text{ dm}^2$}
{$1000\text{ dm}^2$}
{$10\,000\text{ dm}^2$}
\loigiai{
$1\text{ m}^2 = 100\text{ dm}^2$.\\
\textbf{Đáp án của câu hỏi là A}
}
\end{ex}

\begin{ex} %[4D4V18]
$1\text{ cm}^2$ bằng bao nhiêu mi-li-mét vuông?
\choice
{$10\text{ mm}^2$}
{$1000\text{ mm}^2$}
{\True $100\text{ mm}^2$}
{$10000\text{ mm}^2$}
\loigiai{
$1\text{ cm}^2 = 100\text{ mm}^2$.\\
\textbf{Đáp án của câu hỏi là C}
}
\end{ex}

\begin{ex} %[4D4V18]
$1\text{ m}^2$ bằng bao nhiêu mi-li-mét vuông?
\choice
{$1000\text{ mm}^2$}
{\True $1\,000\,000\text{ mm}^2$}
{$100\,000\text{ mm}^2$}
{$10\,000\text{ mm}^2$}
\loigiai{
$1\text{ m}^2 = 100\text{ dm}^2 = 10\,000\text{ cm}^2 = 1\,000\,000\text{ mm}^2$.\\
\textbf{Đáp án của câu hỏi là B}
}
\end{ex}

\begin{ex} %[4D4V18]
$4\text{ m}^2 25\text{ dm}^2$ bằng bao nhiêu đề-xi-mét vuông?
\choice
{$4250\text{ dm}^2$}
{$4025\text{ dm}^2$}
{$42500\text{ dm}^2$}
{\True $425\text{ dm}^2$}
\loigiai{
$4 \times 100 + 25 = 425\text{ dm}^2$.\\
\textbf{Đáp án của câu hỏi là D}
}
\end{ex}

\begin{ex} %[4D4V18]
$300\text{ mm}^2$ bằng bao nhiêu xăng-ti-mét vuông?
\choice
{\True $3\text{ cm}^2$}
{$30\text{ cm}^2$}
{$0{,}3\text{ cm}^2$}
{$300\text{ cm}^2$}
\loigiai{
$300 : 100 = 3\text{ cm}^2$.\\
\textbf{Đáp án của câu hỏi là A}
}
\end{ex}

\begin{ex} %[4D4V18]
Một hình vuông có cạnh dài $10\text{ cm}$. Diện tích của hình vuông đó là
\choice
{$40\text{ cm}^2$}
{$10\text{ dm}^2$}
{\True $1\text{ dm}^2$ ($100\text{ cm}^2$)}
{$100\text{ dm}^2$}
\loigiai{
$10 \times 10 = 100\text{ cm}^2 = 1\text{ dm}^2$.\\
\textbf{Đáp án của câu hỏi là C}
}
\end{ex}

\begin{ex} %[4D4V18]
Một mảnh bìa hình chữ nhật có chiều dài $20\text{ cm}$, chiều rộng $15\text{ cm}$. Diện tích của mảnh bìa đó theo đề-xi-mét vuông là
\choice
{$30\text{ dm}^2$}
{\True $3\text{ dm}^2$}
{$300\text{ dm}^2$}
{$0{,}3\text{ dm}^2$}
\loigiai{
$20 \times 15 = 300\text{ cm}^2 = 3\text{ dm}^2$.\\
\textbf{Đáp án của câu hỏi là B}
}
\end{ex}

\begin{ex} %[4D4V18]
Người ta lát sàn một căn phòng hình chữ nhật dài $6\text{ m}$, rộng $4\text{ m}$ bằng các viên gạch men hình vuông cạnh $4\text{ dm}$. Số viên gạch men cần dùng là
\choice
{$100$ viên}
{$200$ viên}
{$240$ viên}
{\True $150$ viên}
\loigiai{
Diện tích phòng: $6 \times 4 = 24\text{ m}^2 = 2400\text{ dm}^2$.\\
Diện tích 1 viên gạch: $4 \times 4 = 16\text{ dm}^2$.\\
Số viên gạch: $2400 : 16 = 150$ viên.\\
\textbf{Đáp án của câu hỏi là D}
}
\end{ex}

\begin{ex} %[4D4V18]
Điền dấu thích hợp vào chỗ chấm: $5\text{ m}^2 8\text{ dm}^2 \dots 580\text{ dm}^2$.
\choice
{\True $<$}
{$>$}
{$=$}
{Không so sánh được}
\loigiai{
$5\text{ m}^2 8\text{ dm}^2 = 508\text{ dm}^2 < 580\text{ dm}^2$.\\
\textbf{Đáp án của câu hỏi là A}
}
\end{ex}

\begin{ex} %[4D4V18]
Tính giá trị biểu thức: $12\text{ m}^2 - 350\text{ dm}^2 + 50\text{ dm}^2$ theo đơn vị đề-xi-mét vuông.
\choice
{$800\text{ dm}^2$}
{$950\text{ dm}^2$}
{\True $900\text{ dm}^2$}
{$1000\text{ dm}^2$}
\loigiai{
$1200 - 350 + 50 = 900\text{ dm}^2$.\\
\textbf{Đáp án của câu hỏi là C}
}
\end{ex}

\begin{ex} %[4D4V18]
Một cái nhãn vở hình chữ nhật có kích thước $60\text{ mm}$ và $40\text{ mm}$. Diện tích của nhãn vở đó là
\choice
{$240\text{ mm}^2$}
{\True $24\text{ cm}^2$ ($2400\text{ mm}^2$)}
{$240\text{ cm}^2$}
{$200\text{ mm}^2$}
\loigiai{
$60 \times 40 = 2400\text{ mm}^2 = 24\text{ cm}^2$.\\
\textbf{Đáp án của câu hỏi là B}
}
\end{ex}

\begin{ex} %[4D4V18]
Một tấm tôn hình vuông có chu vi $8\text{ m}$. Người ta cắt đi ở 4 góc 4 hình vuông nhỏ bằng nhau, mỗi hình vuông có cạnh $2\text{ dm}$. Diện tích phần tôn còn lại là
\choice
{$384\text{ dm}^2$}
{$396\text{ dm}^2$}
{$360\text{ dm}^2$}
{\True $384\text{ dm}^2$}
\loigiai{
Cạnh tấm tôn: $8 : 4 = 2\text{ m} = 20\text{ dm}$.\\
Diện tích tấm tôn ban đầu: $20 \times 20 = 400\text{ dm}^2$.\\
Diện tích 4 góc cắt: $4 \times (2 \times 2) = 16\text{ dm}^2$.\\
Diện tích còn lại: $400 - 16 = 384\text{ dm}^2$.\\
\textbf{Đáp án của câu hỏi là D}
}
\end{ex}

\textbf{Phần 2. Trắc nghiệm đúng sai}
\setcounterpageref{ex}{0}

\begin{ex} %[4D4V18]
Xét tính đúng sai của các phép đổi đơn vị đo diện tích:
\choiceTF
{\True $3\text{ m}^2 = 300\text{ dm}^2$.}
{$2\text{ dm}^2 5\text{ cm}^2 = 250\text{ cm}^2$.}
{\True $2\text{ dm}^2 5\text{ cm}^2 = 205\text{ cm}^2$.}
{\True $4\text{ cm}^2 = 400\text{ mm}^2$.}
\loigiai{
\begin{itemize}
    \item $3\text{ m}^2 = 300\text{ dm}^2$. Đúng.
    \item $2\text{ dm}^2 5\text{ cm}^2 = 205\text{ cm}^2 \neq 250\text{ cm}^2$. Ý 2 Sai, Ý 3 Đúng.
    \item $4\text{ cm}^2 = 400\text{ mm}^2$. Đúng.
\end{itemize}
\begin{itemize}
    \item a) \textbf{Đúng}.
    \item b) \textbf{Sai}.
    \item c) \textbf{Đúng}.
    \item d) \textbf{Đúng}.
\end{itemize}
}
\end{ex}

\begin{ex} %[4D4V18]
Một mảnh vườn hình chữ nhật có chiều dài $15\text{ m}$, chiều rộng $8\text{ m}$. Xét tính đúng sai:
\choiceTF
{\True Chu vi mảnh vườn là $46\text{ m}$.}
{\True Diện tích mảnh vườn là $120\text{ m}^2$.}
{Đổi diện tích mảnh vườn được $1200\text{ dm}^2$.}
{\True Đổi diện tích mảnh vườn được $12\,000\text{ dm}^2$.}
\loigiai{
\begin{itemize}
    \item Chu vi: $(15 + 8) \times 2 = 46\text{ m}$. Đúng.
    \item Diện tích: $15 \times 8 = 120\text{ m}^2$. Đúng.
    \item $120\text{ m}^2 = 12\,000\text{ dm}^2 \neq 1200\text{ dm}^2$. Ý 3 Sai, Ý 4 Đúng.
\end{itemize}
\begin{itemize}
    \item a) \textbf{Đúng}.
    \item b) \textbf{Đúng}.
    \item c) \textbf{Sai}.
    \item d) \textbf{Đúng}.
\end{itemize}
}
\end{ex}

\textbf{Phần 3. Trả lời ngắn}
\setcounterpageref{ex}{0}

\begin{ex} %[4D4V18]
Đổi $6\text{ m}^2 45\text{ dm}^2$ ra đề-xi-mét vuông.
\shortans{645}
\loigiai{
\textbf{Đáp số:} $645$\\
$645\text{ dm}^2$.
}
\end{ex}

\begin{ex} %[4D4V18]
Đổi $8\text{ cm}^2 30\text{ mm}^2$ ra mi-li-mét vuông.
\shortans{830}
\loigiai{
\textbf{Đáp số:} $830$\\
$830\text{ mm}^2$.
}
\end{ex}

\begin{ex} %[4D4V18]
Một tấm thảm hình vuông có cạnh $30\text{ dm}$. Tính diện tích tấm thảm theo mét vuông.
\shortans{9}
\loigiai{
\textbf{Đáp số:} $9$\\
$30 \times 30 = 900\text{ dm}^2 = 9\text{ m}^2$.
}
\end{ex}

\begin{ex} %[4D4V18]
Một hình chữ nhật có diện tích $450\text{ cm}^2$, chiều rộng $15\text{ cm}$. Tính chiều dài hình chữ nhật theo đề-xi-mét.
\shortans{3}
\loigiai{
\textbf{Đáp số:} $3$\\
Chiều dài: $450 : 15 = 30\text{ cm} = 3\text{ dm}$.
}
\end{ex}

\begin{ex} %[4D4V18]
Cần bao nhiêu viên gạch men hình vuông cạnh $2\text{ dm}$ để lát kín sàn phòng diện tích $20\text{ m}^2$?
\shortans{500}
\loigiai{
\textbf{Đáp số:} $500$\\
$20\text{ m}^2 = 2000\text{ dm}^2 \Rightarrow 2000 : (2 \times 2) = 500$ viên.
}
\end{ex}

\begin{ex} %[4D4V18]
Một hình chữ nhật có chu vi $60\text{ cm}$, chiều dài gấp đôi chiều rộng. Tính diện tích hình chữ nhật đó theo đề-xi-mét vuông (dùng số thập phân hoặc $\text{dm}^2$).
\shortans{2}
\loigiai{
\textbf{Đáp số:} $2$\\
Nửa chu vi: $30\text{ cm}$. Chiều rộng: $10\text{ cm} = 1\text{ dm}$, Chiều dài: $20\text{ cm} = 2\text{ dm}$.\\
Diện tích: $2 \times 1 = 2\text{ dm}^2$.
}
\end{ex}
